% ============================================================
% conclusion.tex  [ENGLISH — RCR submission version]
% Paper: Municipal Public Cleaning Expenditure and Solid Waste
%        Management: Evidence from Peru, 2015–2019
% Author: Dante Barreto Gamarra (UNAS)
% ============================================================

\section{Conclusions}
\label{sec:conclusiones}

This paper estimated the effect of accrued public cleaning expenditure on municipal solid waste management in Peru from 2015 to 2019, using a panel of 9,322 observations across 1,874 districts and 196 provinces. The identification strategy combines the Correlated Random Effects with control function estimator \citep{Wooldridge2015} and the \citet{Mundlak1978} device to correct for expenditure endogeneity without invoking treatment random variation assumptions. The conditional-causal interpretation of the results rests on the CRE assumption: that, controlling for permanent individual heterogeneity through Mundlak means, no heterogeneous unobservable temporal trends persist across municipalities. In the absence of an external instrument --- FONCOMUN integration as an IV remains pending the merge with MEF/SIAF data --- the estimates should be read as conditional effects under this assumption, validated by the \citet{Oster2019} plausibility bounds against static omitted variable bias. The correction is not a technical detail: without it, the TWFE estimator reverses the sign of the effect on final disposal and understates the collection effect by a factor of thirty.

The three main results are precise and statistically robust to multiplicity correction. Expenditure raises daily waste collection by 8.79\% for each 10\% budget increase (elasticity $= +0.879$). It increases the municipal planning index by 0.030 points per 10\% increase ($\hat{\beta} = +0.300$). And it raises by 2.2 percentage points the probability that a municipality disposes its waste in a formal sanitary landfill (AME $= +0.022$). All three primary hypotheses reject the null under Bonferroni-Holm.

These results carry a direct normative implication within the legal framework governing municipal management in Peru. The Organic Law of Municipalities (Law No.\ 27972) establishes solid waste management as an exclusive competence of district local governments. This competence assignment is efficient in theory --- municipalities have the local information to design the service --- but it exposes management outcomes to each district's budget constraint. The results of this paper confirm that the budget constraint is binding: the near-unit elasticity implies that collection capacity scales almost proportionally with expenditure, so the coverage gap between rich and poor municipalities primarily reflects the fiscal gap, not a technical efficiency gap. The state modernization process, embodied in Legislative Decree No.\ 1278 (Solid Waste Integrated Management Law) and Decree No.\ 004-2013-EF introducing performance budgeting, provides the institutional framework for translating this evidence into concrete reforms: linking FONCOMUN distribution to the solid waste management deficit --- measured by the planning index or the open dumpsite rate --- would increase financing precisely where the marginal return on expenditure is highest.

The institutional channel documented in this paper has a second policy implication that goes beyond budget reallocation. The planning index component with the largest marginal effect is PIGARS ($+0.055$ per 10\% expenditure increase), and Legislative Decree No.\ 1278 establishes PIGARS approval as an eligibility condition for MINAM final disposal infrastructure financing. This creates a capacity trap for the poorest municipalities: without expenditure to prepare the PIGARS, they cannot access MINAM financing; without MINAM financing, they cannot build the sanitary landfill that would improve their AME on RS. Municipal expenditure acts as the entry key to the external financing circuit --- not only as a direct service input. A policy covering both constraints simultaneously --- sufficient operating expenditure to plan, and capital co-financing to build --- is more effective than either instrument alone.

Three limitations delimit the scope of the conclusions. The 2015--2019 period predates the disruptive effects of the COVID-19 pandemic on municipal budgets and waste generation patterns; the stability of parameters in the post-2020 period requires verification. The CRE+CF estimator controls for permanent unobservable institutional heterogeneity, but does not purge time-varying confounders without an external instrument; integrating FONCOMUN as an instrumental variable --- pending the merge with MEF/SIAF data --- will strengthen identification in future work. Finally, RENAMU data are statutory declarations without independent physical audit: aspirational reporting bias in planning and final disposal indicators cannot be dismissed.

The body of evidence indicates that the solid waste management gap among Peruvian municipalities is financeable --- not merely technically treatable. Municipalities that today dispose their waste in open dumpsites do not do so because they lack legal incentives (Law No.\ 27972 and Legislative Decree No.\ 1278 establish clear obligations) nor because they are inefficient in using the resources they have. They do so because the resources they have are insufficient to fulfill the obligations the law assigns to them. This is the gap that fiscal policy can close.
